%===============================================================================
\section{Kết luận và Đóng góp khoa học}
%===============================================================================

Quá trình đối chiếu 39 thành phần kiến trúc và 56 bước quy trình khẳng định repository MLOps PaaS đã sở hữu nền tảng vững chắc về quản trị vòng đời mô hình bất biến, hạ tầng huấn luyện đa dạng và kênh thu thập sự kiện suy luận. Đóng góp khoa học cốt lõi của đề tài không phải là việc xây dựng thêm một nền tảng MLOps platform hay gắn thêm trigger tái huấn luyện cơ học, mà là đề xuất và kiểm chứng \textbf{Khung ra quyết định nhận thức suy giảm (Degradation-Aware Decision Framework)} giúp hệ thống tự thích ứng an toàn và tối ưu chi phí. Trong cấu trúc này, repository MLOps PaaS đóng vai trò là \textbf{nền tảng thực nghiệm nguồn mở (Experimental Platform)} phục vụ đánh giá và tái lập kết quả.

Ba đóng góp khoa học chính (Scientific Contributions) bao gồm:
\begin{enumerate}
  \item \textbf{Khung phân loại nguyên nhân và mức độ suy giảm (Degradation Taxonomy \& Severity Model):} Hệ thống hóa 7 nhóm nguyên nhân suy giảm và 5 cấp độ nghiêm trọng ($S_0$--$S_4$), phân tách rành mạch hiện tượng trôi dạt dữ liệu ($P(X)$) khỏi suy giảm hiệu năng thực tế ($P(Y|X)$) và sự cố chất lượng dữ liệu.
  \item \textbf{Chính sách can thiệp thích ứng nhận thức rủi ro (Risk-Aware Adaptive Intervention Policy):} Thuật toán chẩn đoán có thể giải thích được và chính sách điều phối hành động đa nhánh (No-op, Data Engineering, Incremental CT, Full CT, HITL), tích hợp cơ chế watermark xử lý nhãn trễ và quản trị ngân sách tính toán.
  \item \textbf{Giao thức Đánh giá Thực nghiệm Chuẩn hóa (Standardized Benchmark Protocol):} Bộ benchmark toàn diện gồm 12 kịch bản ($S_0$--$S_{11}$) có gắn sẵn nguyên nhân suy giảm thực tế (\codepath{ground\_truth\_cause}), 6 baseline đối chứng ($B_0$--$B_5$) và hệ thống chỉ số đo lường hiệu quả cắt giảm CT thừa ($\text{URR}$), ngăn ngừa bỏ sót suy giảm ($\text{MDR}$) và bảo đảm an toàn tự động hóa.
\end{enumerate}
